\section{Additional Results}

In this section, we provide additional examples generated by all methods compared in all three major experiments we introduced in the main paper, as well as a qualitative illustration of the fidelity-controllability trade-off discussed in Section 4.3.

\paragraph{Generating with COCO Bounding Boxes}
Figure~\ref{fig:coco_extra} shows images generated by all compared methods with COCO bounding boxes. Our proposed CAC is able to enhance consistency between generated images and bounding box information and produce more recognizable objects while maintaining high fidelity.

\paragraph{Generating with Cityscapes Semantic Segmentation Maps}
Figure~\ref{fig:cityscapes_extra} contains additional examples of generated images with Cityscapes semantic segmentation maps. Methods with CAC can generate more coherent and accurate images in comparison to methods without CAC.
\paragraph{Generating with CC-500 Compositional Prompts} 
Figure~\ref{fig:cc500_extra} provide extra samples based on CC-500 captions. CAC improves the compositional generation quality by producing more accurate compositional relationships between objects and attributes.

\paragraph{Ablation Study}
Figure~\ref{fig:ablation_example} is an illustration of the fidelity-controllability trade-off that has been discussed in Section 4.3. Here we also choose MultiDiffusion as the method of controlling self attention in Stable Diffusion and ``MD ratio'' represents the proportion of initial diffusion timesteps that MultiDiffusion is applied to. As we can observe from the figure, without CAC, when very few steps have self attention control, the model cannot generate images that are consistent with the bounding boxes. As the ratio of self attention controlled steps increases, the generation becomes more faithful to the bounding box constraints but loses its fidelity. However, with CAC applied, the model is able to reach a sweet spot where the generated images still appear realistic while maintaining consistency with the bounding box information. This demonstrates the better fidelity-controllability trade-off provided by the CAC application.

\paragraph{Generation with a Variety of User Inputs}
In Figure~\ref{fig:arbitrary} we demonstrate additional examples of using our method to generate images with a variety of user inputs. This showcases the flexibility and effectiveness of CAC in different applications.
\begin{figure}[t]
  \centering
  \includegraphics[width=0.83\textwidth]{img/coco_extra.pdf}
  \caption{Additional examples of generated images based on COCO bounding boxes. ``SD'' denotes Stable Diffusion, ``MD'' denotes MultiDiffusion, and ``Composable'' denotes Composable Diffusion. CAC makes the generated images more recognizable and more accurate to the bounding boxes.}
  \label{fig:coco_extra}
\end{figure}
\begin{figure}[t]
  \centering
  \includegraphics[width=0.85\textwidth]{img/cityscapes_extra.pdf}
  \caption{Additional examples of generated images based on Cityscapes semantic segmentation maps. ``SD'' denotes Stable Diffusion, ``MD'' denotes MultiDiffusion, and ``Composable'' denotes Composable Diffusion. Applying CAC can make the generated images more coherent while being more accurate to the semantic segmentation maps. }
  \label{fig:cityscapes_extra}
\end{figure}
\begin{figure}[t]
  \centering
  \includegraphics[width=\textwidth]{img/cc500_extra.pdf}
  \caption{Additional examples of generated images based on CC-500 captions. ``SD'', ``MD'' and ``Composable'' denote Stable Diffusion, MultiDiffusion and Composable Diffusion. Methods with CAC generate images with more accurate compositional relationships between objects and attributes.}
  \label{fig:cc500_extra}
\end{figure}
\begin{figure}[t]
  \centering
  \includegraphics[width=\textwidth]{img/ablation_example.pdf}
  \caption{Comparison of the fidelity-controllability trade-offs with and without CAC. Sith CAC applied, the model is able to reach a sweet spot where the generated images appear more realistic while maintaining better consistency with the bounding box information. Samples from the same prompts with higher resolution are also shown in Figure~\ref{fig:coco_extra}.}
  \label{fig:ablation_example}
\end{figure}

\begin{figure}[t]
  \centering
  \includegraphics[width=\textwidth]{img/extra_arbitrary.pdf}
  \caption{Examples of generated images with a variety of different types of user inputs and applications.}
  \label{fig:arbitrary}
\end{figure}